% !TEX program = xelatex
\documentclass[11pt,a4paper]{article}
\usepackage{../../../00_common/preamble}

\title{2021年度 (AY2021) 同志社大学大学院 生命医科学研究科\\
医工学コース 入学試験問題「数学」解答例}
\author{}
\date{}

\begin{document}
\maketitle
\vspace{-3em}

%======================================================================
\section*{第1問}

\subsection*{前提整理}

陰関数 $x-y+\log(xy)=0$（$x,y>0$）の $\dfrac{d^2y}{dx^2}$ を求める.
$\log(xy)=\log x+\log y$ と分け,両辺を $x$ で微分する.

\subsection*{計算}

\paragraph{Step 1: 1階導関数.}
$x-y+\log x+\log y=0$ を $x$ で微分すると($y=y(x)$)
\[
1-y'+\frac1x+\frac{y'}{y}=0
\ \Longrightarrow\
\frac{x+1}{x}+y'\,\frac{1-y}{y}=0.
\]
よって
\[
y'=-\frac{x+1}{x}\cdot\frac{y}{1-y}=\frac{(x+1)\,y}{x\,(y-1)}.
\]

\paragraph{Step 2: 2階導関数.}
$y'=\dfrac{(x+1)y}{x(y-1)}$ を商の微分法で $x$ について微分し,現れる $y'$ に
上式を代入して整理する.分子・分母を通分すると
\[
\frac{d^2y}{dx^2}
=-\frac{y\bigl[(x+1)^2+(y-1)^2\bigr]}{x^2\,(y-1)^3}.
\]
\[
\ans{\
\frac{d^2y}{dx^2}=-\frac{y\bigl[(x+1)^2+(y-1)^2\bigr]}{x^2\,(y-1)^3}\ }
\]

検算:曲線上の点 $(x,y)=(2,\ 4.1055\ldots)$ で,有限差分による $y''$ は
$-0.63894\ldots$,上式の値も $-0.63894\ldots$ で一致する.$\checkmark$

%======================================================================
\section*{第2問}

\subsection*{前提整理}

単体 $D:\ x+y+z\le\dfrac\pi2,\ x,y,z\ge0$ 上で
$\displaystyle\iiint_D\cos(x+y+z)\,dV$ を求める.
被積分関数が $s=x+y+z$ のみの関数なので,$s$ 一定の断面積を用いて $1$ 変数積分に直す.

\subsection*{計算}

\paragraph{Step 1: 断面で積分する.}
平面 $x+y+z=s$（$0\le s\le\frac\pi2$）と第1象限の交わりは直角二等辺三角形で,
その面積は $\dfrac{s^2}{2}$.よって
\[
\iiint_D\cos(x+y+z)\,dV=\int_0^{\pi/2}\cos s\cdot\frac{s^2}{2}\,ds.
\]

\paragraph{Step 2: 部分積分する.}
$\displaystyle\int s^2\cos s\,ds=s^2\sin s+2s\cos s-2\sin s$ なので
\[
\int_0^{\pi/2}s^2\cos s\,ds
=\Bigl[s^2\sin s+2s\cos s-2\sin s\Bigr]_0^{\pi/2}
=\frac{\pi^2}{4}-2.
\]
したがって
\[
\ans{\
\iiint_D\cos(x+y+z)\,dx\,dy\,dz=\frac12\Bigl(\frac{\pi^2}{4}-2\Bigr)=\frac{\pi^2}{8}-1\ }
\]

検算:$\dfrac{\pi^2}{8}-1=0.2337\ldots>0$.被積分関数は $D$ 上で
$0\le x+y+z\le\frac\pi2$ ゆえ $\cos\ge0$ で正,符号と整合する.$\checkmark$

%======================================================================
\section*{第3問}

\subsection*{前提整理}

\[
A=\begin{pmatrix}a&1&0&0\\0&1&0&0\\b&0&b&0\\0&0&1&c\end{pmatrix}
\qquad(a,b,c\neq0)
\]
の行列式・固有値・固有ベクトルを求める.

\subsection*{計算}

\paragraph{Step 1: 行列式.}
第4列 $(0,0,0,c)^\top$ で展開し,続く $3$ 次行列式を第3列で展開すると
\[
|A|=c\begin{vmatrix}a&1&0\\0&1&0\\b&0&b\end{vmatrix}
=c\cdot b\begin{vmatrix}a&1\\0&1\end{vmatrix}=abc.
\]
よって $\ans{|A|=abc}$.

\paragraph{Step 2: 固有値.}
同じ展開を $A-\lambda I$ に施すと
\[
\det(A-\lambda I)=(c-\lambda)(b-\lambda)(a-\lambda)(1-\lambda),
\]
したがって固有値は $\lambda=a,\ 1,\ b,\ c$.（以下 $a,b,c,1$ は相異なるとする.）

\paragraph{Step 3: 固有ベクトル.}
各 $(A-\lambda I)\vv=\vzero$ を解くと
\[
\ans{\
\begin{aligned}
&\lambda=c:\ \vv=\begin{psmallmatrix}0\\0\\0\\1\end{psmallmatrix},\quad
\lambda=b:\ \vv=\begin{psmallmatrix}0\\0\\b-c\\1\end{psmallmatrix},\\[4pt]
&\lambda=a:\ \vv=\begin{psmallmatrix}(a-b)(c-a)\\0\\b(c-a)\\-b\end{psmallmatrix},\quad
\lambda=1:\ \vv=\begin{psmallmatrix}(b-1)(c-1)\\-(a-1)(b-1)(c-1)\\-b(c-1)\\b\end{psmallmatrix}
\end{aligned}\ (\text{各定数倍})\ }
\]

検算:固有値の積 $a\cdot1\cdot b\cdot c=abc=|A|$,和 $a+1+b+c=\operatorname{tr}A$ と一致.
また $\lambda=b$ では $A\vv$ の第3成分 $b\cdot0+b(b-c)+0=b(b-c)=b\cdot(b-c)$,
$\vv$ の第3成分 $b-c$ の $b$ 倍に等しく整合.$\checkmark$

%======================================================================
\section*{第4問}

\subsection*{前提整理}

\[
W_1:\ \begin{cases}x-2y+w=0\\y+z+w=0\end{cases}\qquad
W_2:\ x-y-z=0
\]
（いずれも $\R^4$ の部分空間）について,$W_1$ の基底・次元と $\dim(W_1+W_2)$ を求める.

\subsection*{計算}

\paragraph{Step 1: $W_1$ の基底と次元.}
$2$ 本の方程式は1次独立なので $\dim W_1=4-2=2$.
$x=2y-w,\ z=-y-w$ と表し,$(y,w)=(1,0),(0,1)$ をとると基底
\[
\ans{\
\dim W_1=2,\quad
\text{基底}\ \Bigl\{(2,1,-1,0)^\top,\ (-1,0,-1,1)^\top\Bigr\}\ }
\]

\paragraph{Step 2: $\dim(W_1+W_2)$.}
$W_2$ は $1$ 本の方程式なので $\dim W_2=3$.
$W_1\cap W_2$ は $3$ 本の方程式
\[
\begin{pmatrix}1&-2&0&1\\0&1&1&1\\1&-1&-1&0\end{pmatrix}
\]
で定まり,この係数行列の階数は $3$（行簡約で確認）だから $\dim(W_1\cap W_2)=4-3=1$.
次元公式より
\[
\dim(W_1+W_2)=\dim W_1+\dim W_2-\dim(W_1\cap W_2)=2+3-1=4.
\]
\[
\ans{\ \dim(W_1+W_2)=4\ (\text{すなわち}\ W_1+W_2=\R^4)\ }
\]

検算:基底ベクトル $(2,1,-1,0)^\top$ は $x-2y+w=2-2+0=0$,$y+z+w=1-1+0=0$ を
満たし $W_1$ に属する.$\dim(W_1+W_2)=4$ は $\R^4$ 全体と一致し,
$2+3=5>4$ から共通部分が $1$ 次元であることとも整合.$\checkmark$

\end{document}
